\section{Weekly Summary}

\begin{tcolorbox}[colback=yellow!10,colframe=black!50,title=AI-Generated Content]
\noindent The summary below was written by Claude (Anthropic) from that week's regression outputs and series descriptions. It has not been reviewed or fact-checked by a human, and should be read as an automated, informal recap -- not analysis.
\end{tcolorbox}

Welcome back to the weekly parade of numbers that have no business being in the same room together, let alone the same scatterplot. This week our little machine reached into the FRED grab bag five times and pulled out such crowd-pleasers as discontinued small-business loan risk percentages, aggregate weekly hours in durable goods manufacturing, imports of non-manufactured commodities to the U.S. Virgin Islands, and -- my personal favorite -- grocery inventories versus Southeastern social assistance GDP. Truly the questions that keep economists up at night, if by "economists" you mean "a script running unsupervised at 3 a.m."

We opened Monday with the discontinued SBA loan series versus Ft. Lauderdale house prices, and I regret to inform the tabloids that there is nothing here. The raw regression shrugged (p around 0.33), and the detrended version shrugged even harder. A rare case of the machine correctly identifying that two random things are, in fact, unrelated. Tuesday gave us the week's most instructive little morality play: durable-goods weekly hours versus foreign official demand deposits looked statistically "significant" in the raw run (p about 0.006, gasp, clutch pearls), and then the instant we stripped out each series' own drifting-over-time habit, the whole thing evaporated to a p-value of 0.82. That, dear reader, is the textbook trend mirage this entire project exists to point and laugh at: two things wandering the same general direction over the years, briefly mistaken for friends.

Wednesday flipped the script in the funniest way. Virgin Islands imports versus Dutch services business confidence had absolutely nothing going on in the raw regression (p 0.59), but the detrended version squeaked into "significance" (p about 0.014) -- with an R-squared of a heart-stopping 0.027, meaning it explains roughly nothing while technically clearing the bar. Thursday's grocery inventories versus social assistance GDP did something similar, coming alive only after detrending (p 0.037, R-squared a positively electric 0.16). I'd theorize a deep macroeconomic linkage between canned goods and the Southeast's welfare economy, but I'd also like to keep my dignity, so I won't.

The genuine star of the week is Friday: New England premium gasoline prices versus net sales of "all other professional and technical services except legal." The raw regression came in swinging with an R-squared of 0.52 and a p-value with more zeros than my patience, and -- plot twist -- it survived detrending too (p about 0.034), which is the rare case where I'm contractually obligated to say "huh, that's actually kind of interesting." Maybe pricier gas and busier consulting-adjacent firms both ride the same broader economic tide; maybe it's a coincidence wearing a very convincing costume. Either way, standard disclaimer, shouted from the rooftops: these are randomly paired series, no theory, no causation, no rigor beyond "the line went diagonal." Spurious correlation is the house special here, and this week the kitchen was cooking as usual.
